\begin{figure}[htbp]
\centering
% frame common to all graphs: cut, Wilson line (arrow at position #1), operator vertices, external labels
\newcommand{\toneW}[1]{%
  \draw[cut] (0,-0.22) -- (0,2.6);
  \draw[wilson=#1] (-0.75,1.9) -- (0.75,1.9);
  \node[opvertex] at (-0.75,1.9) {}; \node[opvertex] at (0.75,1.9) {};
  \node[below] at (-0.75,0) {$p$}; \node[below] at (0.75,0) {$p$};}
\begin{subfigure}[b]{0.24\textwidth}\centering
\begin{tikzpicture}[fd]
  \draw[quark] (-0.75,0) -- (-0.75,0.9);
  \draw[quark] (-0.75,0.9) -- (-0.75,1.9);
  \draw[quark] (0.75,1.9) -- (0.75,0.9);
  \draw[quark] (0.75,0.9) -- (0.75,0);
  \draw[gluon] (-0.75,0.9) -- (0.75,0.9);
  \toneW{0.25}
  \node[vertex] at (-0.75,0.9) {}; \node[vertex] at (0.75,0.9) {};
  \path[mom'={$l$},mom pos=0.76,mom length=0.42cm,mom sep=0.27cm] (-0.75,0.9) -- (0.75,0.9);
  \path[mom={$k$},mom length=0.45cm] (-0.75,0.9) -- (-0.75,1.9);
\end{tikzpicture}
\caption{实胶子发射}
\end{subfigure}\hspace{2em}
\begin{subfigure}[b]{0.56\textwidth}\centering
\begin{tikzpicture}[fd]
  \draw[quark] (-0.75,0) -- (-0.75,0.8);
  \draw[quark] (-0.75,0.8) -- (-0.75,1.9);
  \draw[quark] (0.75,1.9) -- (0.75,0);
  \draw[gluon] (-0.75,0.8) to[out=10,in=-100] (0.4,1.9);
  \toneW{0.25}
  \node[vertex] at (-0.75,0.8) {}; \node[vertex] at (0.4,1.9) {};
\end{tikzpicture}\hspace{1.5em}
\begin{tikzpicture}[fd]
  \draw[quark] (-0.75,0) -- (-0.75,1.9);
  \draw[quark] (0.75,1.9) -- (0.75,0.8);
  \draw[quark] (0.75,0.8) -- (0.75,0);
  \draw[gluon] (-0.4,1.9) to[out=-80,in=170] (0.75,0.8);
  \toneW{0.75}
  \node[vertex] at (0.75,0.8) {}; \node[vertex] at (-0.4,1.9) {};
\end{tikzpicture}
\caption{夸克线与另一侧 Wilson~线间的实胶子（及镜像）}
\end{subfigure}

\bigskip
\begin{subfigure}[b]{0.68\textwidth}\centering
\begin{tikzpicture}[fd]
  \draw[quark] (-0.75,0) -- (-0.75,0.8);
  \draw[quark] (-0.75,0.8) -- (-0.75,1.9);
  \draw[quark] (0.75,1.9) -- (0.75,0);
  \draw[gluon] (-0.75,0.8) to[out=0,in=-90] (-0.25,1.9);
  \toneW{0.75}
  \node[vertex] at (-0.75,0.8) {}; \node[vertex] at (-0.25,1.9) {};
\end{tikzpicture}\hspace{0.9em}
\begin{tikzpicture}[fd]
  \draw[quark] (-0.75,0) -- (-0.75,1.9);
  \draw[quark] (0.75,1.9) -- (0.75,0.8);
  \draw[quark] (0.75,0.8) -- (0.75,0);
  \draw[gluon] (0.25,1.9) to[out=-90,in=180] (0.75,0.8);
  \toneW{0.25}
  \node[vertex] at (0.75,0.8) {}; \node[vertex] at (0.25,1.9) {};
\end{tikzpicture}\hspace{0.9em}
\begin{tikzpicture}[fd]
  \draw[quark] (-0.75,0) -- (-0.75,0.45);
  \draw[quark] (-0.75,0.45) -- (-0.75,1.3);
  \draw[quark] (-0.75,1.3) -- (-0.75,1.9);
  \draw[quark] (0.75,1.9) -- (0.75,0);
  \draw[gluon] (-0.75,0.45) arc[start angle=-90,end angle=-270,radius=0.425];
  \toneW{0.25}
  \node[vertex] at (-0.75,0.45) {}; \node[vertex] at (-0.75,1.3) {};
\end{tikzpicture}\hspace{0.9em}
\begin{tikzpicture}[fd]
  \draw[quark] (-0.75,0) -- (-0.75,1.9);
  \draw[quark] (0.75,1.9) -- (0.75,1.3);
  \draw[quark] (0.75,1.3) -- (0.75,0.45);
  \draw[quark] (0.75,0.45) -- (0.75,0);
  \draw[gluon] (0.75,1.3) arc[start angle=90,end angle=-90,radius=0.425];
  \toneW{0.25}
  \node[vertex] at (0.75,0.45) {}; \node[vertex] at (0.75,1.3) {};
\end{tikzpicture}
\caption{不穿过割线的虚图（两两互为镜像）}
\end{subfigure}\hfill
\begin{subfigure}[b]{0.28\textwidth}\centering
\begin{tikzpicture}[fd]
  \draw[quark] (-0.75,0) -- (-0.75,1.9);
  \draw[quark] (0.75,1.9) -- (0.75,0);
  \draw[gluon] (-0.5,1.9) arc[start angle=180,end angle=0,radius=0.5];
  \toneW{0.36}
  \node[vertex] at (-0.5,1.9) {}; \node[vertex] at (0.5,1.9) {};
  \node[font=\footnotesize] at (0,-0.68) {$\propto n^2=0$};
\end{tikzpicture}
\caption{Wilson~线之间的胶子}
\end{subfigure}
\caption{Feynman 规范下夸克中夸克分布 $f_{q/q}(\xi)$ 的 $\order{\as}$ 割图（靶为在壳夸克 $p$，$k^+=\xi p^+$；记号同\cref{fig:pdf-operator}）。
(a) 实胶子发射（梯形图），给出 $P_{qq}$ 中的 $C_F(1-\xi)$；(b) 实胶子由夸克线连到割线另一侧的 Wilson 线，及其镜像；
(c) 不穿过割线的虚图：夸克线—Wilson 线顶点修正（前两图）与外线自能（后两图）。
(b) 与 (c) 中夸克线—Wilson 线图在 $\xi\to1$ 的端点奇异性只在实、虚之和中相消，合成 $C_F[2\xi/(1-\xi)]_+$；
外线自能的紫外发散（夸克场重整化）贡献 $-\tfrac12C_F\delta(1-\xi)$；总和为 $P_{qq}=C_F[(1+\xi^2)/(1-\xi)]_+$。
(d) 连接割线两侧 Wilson 线的胶子正比于 $n^2=0$，在 Feynman 规范中为零（同侧的 Wilson 线自能亦然）。}
\label{fig:pdf-oneloop}
\end{figure}
